\scanpage{25}
\handwritten{SGA 4 \uline{IX}\quad p. 42}
\margin{\handwritten{nous prouverons en même temps que $\tilde{X}_{\mathrm{et}}$ satisfait aux conditions équivalentes de \uline{VI} 1.2\unsure{).|].}}}

\handwritten{(*) (Note pour la page 8.) Prouvons que la notion de constructibilité 2.3 pour un faisceau d'ensembles $F$ coïncide avec celle de \uline{VI} 1.9.3. \deleted{Soit} La \deleted{\unsure{en effet $C$ le cas}}\inserted{question étant} locale sur $X$, on peut supposer $X$ affine \inserted{donc ``cohérent'', i.e.} quasi-compact et quasi-séparé. \deleted{Soit $C$ la catégorie \unsure{Soit $(C')$ le sous}\illegible{1}} \deleted{\unsure{Il|Je} faut alors prouver que $F$ est constructible} Soit $C$ (resp. $C'$) la sous-catégorie \deleted{\unsure{pleine}}\inserted{strictement} pleine de $\tilde{X}_{\mathrm{et}}$ formée des $F$ constructibles au sens de 2.3, \deleted{\illegible{1}} \inserted{(resp. au sens de \uline{VI} 1.9.3), il} faut prouver \deleted{que $C$ s'identifie à la catégorie des objets \uline{cohérents} de} \deleted{\unsure{$\tilde{X}_{\mathrm{et}}$}} que $C=C'$;}

\handwritten{\deleted{En vertu de \uline{VII} 5.6 a), \deleted{l'objet final $X$ de} $\tilde{X}_{\mathrm{et}}$ \deleted{est cohérent, donc} $C'$ est formé des objets \uline{cohérents} \deleted{\unsure{\uline{VI}}} de $\tilde{X}_{\mathrm{et}}$ (2.5). En vertu de \uline{VI} 1.24, \deleted{\unsure{et}} et compte tenu du fait que $C$ est génératrice dans $\tilde{X}_{\mathrm{et}}$ (2.7.2) et stable par limites inductives finies (2.\unsure{6|8} (i)), il suffit de prouver que $C\subset C'$, Mais il résulte de 2.7 que tout $F\in\mathrm{ob}\,C$ est \uline{quasi-compact}, puisque un \deleted{\unsure{$G\in\mathrm{ob}\,X_{\mathrm{et}}$}}\inserted{$\unsure{G|C}\in X'_{\mathrm{et}}$} \deleted{\unsure{schéma $X'$}} qui est quasi-compact et quasi-séparé est quasi-compact comme objet de $\tilde{X}_{\mathrm{et}}$ \deleted{(5.6.1)} (\uline{VII} 5.6.1). \deleted{D'autre part pour \unsure{voir que} $\unsure{F|F'}\unsure{\in|\to}\unsure{\mathrm{ob}\,X'|X'}$ la diagonale $F\longrightarrow F\mathbin{\times}_{X}F$ est quasi-compacte, il suffit de voir que \unsure{pour tout} \illegible{2} \unsure{et que}\illegible{2} \unsure{ce fait que $C$ est stable par produits}} Comme $C$ est stable par \deleted{limites projectives}\inserted{produits} finis\inserted{(2.6 (ii))}, l'inclusion $C\subset C'$ est donc un cas particulier de 2.1.}}
\margin{\handwritten{\deleted{\unsure{pourquoi|pour} \unsure{préciser}}}}

\handwritten{En vertu de \uline{VI} 2.9, b), il suffit de voir que $C$ est stable pour limites inductives finies et limites projectives finies, qu'elle est génératrice, et formée d'objets quasi-compacts. \deleted{\illegible{1}} Compte tenu de 2.6 (i) et (2.7.2), il reste à prouver que tout $F\in\mathrm{ob}\,C$ est quasi-compact. Or \deleted{cela résulte}\inserted{en vertu} de 2.7 \deleted{\unsure{et} \uline{VI} 1.3}, \deleted{puisque}\inserted{il existe} un $G\in\mathrm{ob}\,X_{\mathrm{et}}$ \deleted{\illegible{1}} qui est \inserted{un schéma} quasi-compact et quasi-séparé \deleted{\unsure{et est}} et un épimorphisme $G\longrightarrow F$. Or $G$ \deleted{\unsure{comme on}} est un objet quasi-compact de $\tilde{X}_{\mathrm{et}}$ (\uline{VII} 5.6.1), donc $F$ est quasi-compact (\uline{VI} 1.3), qfd.}

\handwritten{Bien entendu, il y aurait \inserted{eu} lieu de refondre les §§ 1,2 du présent exposé, \inserted{notamment} en les rattachant \deleted{\unsure{explicitement}} aux notions générales de \deleted{l'exposé} finitude de \uline{VI}. Nous laissons ce soin \deleted{\unsure{complémentaire}} de rédaction aux générations futures.}
% archive sheet number: 24
